\documentclass{article}
\usepackage{/globals/de}
\usepackage{/globals/anki}
\ankiprefix{gesetzgebungsverfahren}
\ankitopic{Gesetzgebungsverfahren}
\ankishow
\begin{document}
\section{Das Gesetzgebungsverfahren} 
Das Gesetzgebungsverfahren kann in drei Teile aufgeteilt werden; die Initiativphase, die Beschlussphase und die Ausfertigungsphase.
\ankinote{phasen}{Phasen}{Initiativ, Beschluss, Ausfertigung}

\subsection{Initiation}
Die Bundesregierung, der Bundestag und der Bundesrat können Gesetzesvorlagen, Gesetzestexte mit Begründungen, einbringen.
\ankinote{einbringen}{Einbringen}{BReg, BT, BR Gesetzesvorlagen (Texte+Begründung)}

Wenn die Bundesregierung oder der Bundesrat eine Vorlage einbringen muss das jeweils andere Organ Stellung nehmen, damit Streitpunkte früh geklärt werden können. Dafür wird es in einem Übersendungsschreiben übersendet. Zudem nimmt der Nationale Normenkontrollenrat Stellung, zu den Folgekosten.
\ankinote{übersenden}{Übersenden}{Übersendungsschreiben, BReg+BR jeweils anderes Organ Stellung}
\ankinote{stellungnahme}{Stellungnahme}{BReg+BR, NKR Nationaler Normenkontrollenrat}

\subsection{Beschluss}
Gesetze werden in drei Lesungen beschlossen.
\begin{description}
 \item[Lesung I] Der Gesetzesentwurf wird erklärt und an die Fachausschüsse überwiesen.
 \item[Ausschussberatungen] Die Ausschüsse beraten sich, es wird eine Beschlussempfehlung für das Plenum erstellt.
 \item[Lesung II] Die Beschlussempfehlung wird beraten, es kann Änderungsanträge geben.
 \item[Lesung III] Es kann Änderungsanträge geben, zu Punkten die in der zweiten Lesung geändert wurden. 
 \item[Schlussbestimmung] Finale Abstimmung über den Entwurf.
\end{description}
\ankinote{beschluss-phasen}{Beschluss Phasen}{Lesung I, Ausschussberatung, Lesung II, Lesung III, Schlussbestimmung}
\ankinote{les-i}{Lesung I}{Erklärt, Ausschüsse überwiesen}
\ankinote{ausschuss}{Ausschussberatung}{Ausschüsse beraten sich, Beschlussempfehlung}
\ankinote{les-ii}{Lesung II}{Beschlussempfehlung beraten, Änderungsanträge}
\ankinote{les-iii}{Lesung III}{Änderungsanträge zu Punkten von Lesung II}
\ankinote{schlussbestimmung}{Schlussbestimmung}{Finale Abstimmung}

Der genaue Ablauf in Bezug auf die Zusammenarbeit mit dem Bundesrat hängt dabei von der Gesetzesart ab.
\begin{description}
 \item[Zustimmungsgesetze] Gesetze zu bestimmten Materien sind Zustimmungsgesetze; sie können nur mit der Zustimmung vom Bundesrat entschieden werden.
 \item[Einspruchsgesetze] Alle anderen Gesetze sind Einspruchsgesetze; der Bundesrat kann nur Einspruch einlegen, dafür muss der Bundesrat vorher den Vermittlungsausschuss angerufen haben. Auch wenn er Einspruch einlegt kann der Bundestag mit einer größeren Mehrheit diesen zurückweisen.
\end{description}
\ankinote{zustimm}{Zustimmungsgesetze}{Bestimmter Materien; nur mit Zustimmung Bundesrat}
\ankinote{einspruch}{Einspruchsgesetze}{Andere Gesetze; VA, Einspruch; Überstimmung mit größerer Mehrheit}

Wird der Vermittlungsausschuss angerufen, so soll er einen Lösungsvorschlag dem sowohl der Bundestag als auch der Bundesrat zustimmen finden. Dieser muss sich im Rahmen der bisherigen Stellungnahmen und Anträgen befinden. Kommt es zu einem Lösungsvorschlag, so müssen Bundestag und Bundesrat diesem erneut zustimmen.
\ankinote{va}{Vermittlungsausschuss}{Lösungsvorschlag im Rahmen der bisherigen Anträge, neue Zustimmung}

\subsection{Ausfertigung}
Das Gesetz wird ausgefertigt; der Bundespräsident unterzeichnet es, dabei kann er es formell und umstrittenerweise auch materiell überprüfen.

Danach wird es verkündet, es wird im Bundesgesetzblatt veröffentlicht. Damit tritt es in Kraft.
\ankinote{ausfertigung}{Ausferting Phasen}{BuPrä, Veröffentlichung}
\ankinote{buprä}{BuPrä}{Unterzeichnet, formelle und materielle Prüfung}
\ankinote{verkündung}{Verkündung}{Bundesgesetzblatt veröffentlicht, tritt in Kraft}
\end{document}